\section{Sifat Fungsi Hash Kriptografis}

Fungsi hash memetakan input dengan panjang sembarang ke output dengan panjang tetap (misalnya 256 bit) \cite{stallings2017}. Konstruksi Merkle-Damgård menjadi fondasi bagi MD5, SHA-1, SHA-2 \cite{menezes1996}.

Fungsi hash kriptografis harus memenuhi:
\begin{itemize}
  \item \textbf{Preimage resistance}: Dari $h(x)$, sulit mencari $x$
  \item \textbf{Second preimage resistance}: Dari $x$, sulit mencari $x' \neq x$ dengan $h(x') = h(x)$
  \item \textbf{Collision resistance}: Sulit mencari pasangan $(x, x')$ dengan $x \neq x'$ dan $h(x) = h(x')$
\end{itemize}

\textbf{Penggunaan}: Integritas data, penyimpanan password (hash + salt), tanda tangan digital (hash pesan lalu tanda tangan hash) \cite{stallings2017}.

\subsection{Konstruksi Merkle-Damgård}

Konstruksi Merkle-Damgård menggunakan fungsi kompresi collision-resistant dan iterasi blok demi blok \cite{menezes1996}. Dasar bagi MD5, SHA-1, SHA-2. SHA-3 menggunakan konstruksi sponge yang berbeda \cite{stallings2017}.

\subsection{Avalanche Effect}

Perubahan kecil pada input menyebabkan perubahan besar pada output. Setiap bit input mempengaruhi banyak bit output \cite{stallings2017}.

\subsection{Standard Hash Modern}

SHA-2 (SHA-224, SHA-256, SHA-384, SHA-512): standar NIST, aman \cite{stallings2017}. SHA-3 (Keccak): standar alternatif dengan konstruksi sponge; tahan terhadap serangan yang menargetkan Merkle-Damgård \cite{menezes1996}.
